% Level 3, Week 3 -- Families of Objects (student workbook)
\chapterweek{3}{Families of Objects}{Child classes inherit from a parent and change only what is different; bigger objects are built from smaller ones.}

\section{What you will learn}
\begin{checklist}
  \item Make a child class that inherits from a parent class.
  \item Override a method, and use \texttt{super()} to reuse the parent's code.
  \item Call the same method on different objects and get different results (polymorphism).
  \item Build an object that holds a list of other objects (composition).
\end{checklist}

\section{The big idea: inheritance}
A \textbf{child class} gets everything from its \textbf{parent class} for free. Then it changes only what is different. Write the parent in brackets:

\begin{center}
\begin{tikzpicture}[
  cls/.style={draw=black,line width=0.9pt,rounded corners=3pt,align=center,font=\small,minimum width=3.4cm,minimum height=1.25cm,inner sep=4pt}]
  \node[cls,fill=kcblue!15] (enemy) at (0,0) {\textbf{Enemy} (parent)\\\texttt{name, x, y, steps}\\\texttt{move()}: 1 step right};
  \node[cls,fill=kcyellow!35] (guard) at (-5.2,-2.8) {\textbf{Guard}\\\texttt{move()}: stays still};
  \node[cls,fill=kcgreen!20] (zig) at (0,-2.8) {\textbf{Zigzagger}\\\texttt{move()}: parent's move,\\then up or down};
  \node[cls,fill=kcorange!25] (spr) at (5.2,-2.8) {\textbf{Sprinter}\\extra: \texttt{speed}\\\texttt{move()}: \texttt{speed} steps};
  \foreach \c in {guard,zig,spr} \draw[-{Triangle[open,length=3mm,width=3mm]},line width=0.9pt] (\c.north) -- (enemy.south);
  \node[font=\small,align=center] at (3.6,-0.9) {arrow = \enquote{is a}};
\end{tikzpicture}
\end{center}

\begin{pycode}[A child class]
class Sprinter(Enemy):                  # a Sprinter IS AN Enemy
    def __init__(self, name, x=0, y=0, speed=3):
        super().__init__(name, x, y)    # parent sets name, x, y
        self.speed = speed              # then the extra attribute

    def move(self):                     # override the parent's move
        self.x = self.x + self.speed
        self.steps = self.steps + 1
\end{pycode}

\begin{tip}[super() = \enquote{do the parent's part}]
\texttt{super().\_\_init\_\_(name, x, y)} runs the parent's setup. \texttt{super().move()} runs the parent's \texttt{move}, so the child can add something on top. If a child has its own \texttt{\_\_init\_\_} and forgets \texttt{super().\_\_init\_\_}, its parent's attributes are missing and you get an \texttt{AttributeError}.
\end{tip}

\section{Polymorphism: one call, many behaviors}
\begin{pycode}
enemies = [Enemy("Blob"), Guard("Rocky"), Sprinter("Zoom")]
for enemy in enemies:
    enemy.move()        # each one moves its OWN way
\end{pycode}
The loop does not need to know what kind each enemy is. \emph{Polymorphism} means \enquote{many shapes}: the same message, \texttt{move()}, gives different results.

\section{Composition: \enquote{has a}}
Not everything is a family. A Game is not a kind of Player, but a Game \textbf{has} Players. When an object keeps other objects in an attribute (often a list), that is \textbf{composition}.
\begin{center}
\begin{tabular}{@{}p{6.5cm}p{6.5cm}@{}}
\toprule
\textbf{Is a \ding{213} inheritance} & \textbf{Has a \ding{213} composition} \\ \midrule
A Sprinter is an Enemy. & A Game has Players. \\
A Star is a Shape. & A Party has Heroes. \\
\texttt{class Star(Shape):} & \texttt{self.players = []} \\
\bottomrule
\end{tabular}
\end{center}

\begin{exercise}[Enemy Family]
Folder: \texttt{week03/ex-enemy-family}. Finish \texttt{Zigzagger.move} (use \texttt{super().move()}), \texttt{Sprinter.\_\_init\_\_} (use \texttt{super().\_\_init\_\_}), \texttt{Sprinter.move} and \texttt{move\_all}.\\
\textbf{Predict:} a Zigzagger starts at (0, 0). Where is it after 3 moves? \rule{3cm}{0.4pt}
\end{exercise}

\begin{exercise}[Shape Family]
Folder: \texttt{week03/ex-shape-family}. \texttt{Shape} and \texttt{Square} are done. Make \texttt{Triangle} (3 sides, \texttt{left(120)}) and \texttt{Star} (5 sides, \texttt{right(144)}) override \texttt{draw}. Fix \texttt{Polygon.\_\_init\_\_} with \texttt{super()}. Run it: one call to \texttt{draw\_all} paints a whole grid of different shapes.
\end{exercise}

\begin{exercise}[Game Roster (if you finish early)]
Folder: \texttt{week03/ex-game-roster}. A \texttt{Game} \emph{has} a list of \texttt{Player} objects. Finish \texttt{add\_player}, \texttt{give\_points}, \texttt{leader} and \texttt{still\_playing}.
\end{exercise}

\begin{vocab}
  \vocabword{inheritance}{A child class gets all the attributes and methods of its parent class.}
  \vocabword{parent class / child class}{The general class and the more special class built from it. Also called superclass and subclass.}
  \vocabword{override}{Write a method in the child with the same name, to replace the parent's version.}
  \vocabword{super()}{Inside a child class: \enquote{my parent}. Use it to run the parent's version of a method.}
  \vocabword{polymorphism}{The same method call does different things for different classes.}
  \vocabword{composition}{An object that holds other objects, like a Game that has a list of Players.}
\end{vocab}

\begin{reflection}
Draw a small family tree for a game you would like to make: one parent and two children. What does each child do differently?
\answerspace{3cm}
Write one \enquote{is a} sentence and one \enquote{has a} sentence about your game:
\answerlines{2}
\end{reflection}
